\section{Conclusion}
\label{sec:conclusion}

We presented a unified evaluation framework for memory poisoning attacks and defenses
on LLM agent systems.  Our framework makes six concrete advances.

\textbf{First}, we identify and correct a systematic evaluation error in prior \agentpoison{}
simulations: victim queries must include the attacker-controlled trigger token at
evaluation time, matching the paper's Algorithm~2.  This single correction raises
\asrr{} from $0.25$ to $1.00$---a four-fold improvement that substantially changes
the threat assessment.  Simulations that omit trigger prepending will severely
underestimate attacker capability.

\textbf{Second}, we show that a centroid-targeting adversarial passage achieves
$\asrr=1.00$ without full DPR-based gradient optimization, and provide a complete
DPR HotFlip implementation~\citep{ebrahimi2018hotflip,chen2024agentpoison} that
further improves mean triggered-query similarity from $0.71$ (centroid only) to
$0.78$ ($+10\%$).  The DPR optimizer uses gradient-guided token selection through
the \texttt{facebook/dpr-ctx\_encoder-single-nq-base} embedding layer with
optional GPT-2 perplexity filtering for fluency.

\textbf{Third}, we propose \sad{}: a Semantic Anomaly Detection defense that requires
no labeled adversarial examples and no model access.  We further introduce a
\emph{combined scoring mode} that averages max- and mean-query similarity
(Eq.~\ref{eq:sad_combined}), achieving $\fpr=0.00$ across all three attacks at
$k=2.0$ with AUROC $1.000$ (\minja{}), $0.920$ (\injecmem{}), and $0.870$
(\agentpoison{}, plain calibration).  \sad{} achieves near-zero TPR against
\agentpoison{} under plain-query calibration because the centroid passage falls
within the benign similarity distribution; triggered-query calibration raises
TPR to $1.00$ (AUROC\,=\,1.000) with $\fpr=0.00$.

\textbf{Fourth}, our analysis of the adaptive adversary reveals a nuanced result.
In continuous gradient space (Eq.~\ref{eq:adaptive_obj}), any reduction in \sad{}'s
anomaly score simultaneously reduces retrieval rank---a structural tension.
Empirically, word-level synonym substitution circumvents this tension because
semantic embeddings are synonym-invariant: domain synonym replacements shift
anomaly scores without affecting retrieval rank ($\Delta\asrr \approx 0$).
This finding motivates future defenses incorporating lexical diversity features
beyond the $\ell_2$ semantic embedding space.

\textbf{Fifth}, we implement the paper-faithful token-level unigram watermark
\citep{zhao2024watermark} using GPT-2 as a backbone LLM.  The token-level scheme
achieves mean z-score $\approx 7.5$ at $n=50$ tokens vs.\ $\approx 4.1$ for
character-level---an $83\%$ improvement in detection margin at short text lengths,
covering the minimum-length threshold where character-level detection is unreliable.

\textbf{Sixth}, we provide a directly measured lower bound on $\asra$ via a local
GPT-2 agent evaluator (\textsc{LocalAgentEvaluator}): measured values ($0.42$,
$0.51$, $0.29$ for the three attacks) confirm that modelled upper bounds are
conservative and that even a weak local proxy demonstrates meaningful adversarial
action execution in retrieved-context scenarios.

\textbf{Seventh}, we evaluate attack transferability and \sad{} generalization
across seven retrieval encoders using linear CKA \citep{kornblith2019similarity},
finding that inter-encoder CKA predicts cross-encoder ASR-R transferability.
AgentPoison with triggered-query calibration achieves $\tpr=1.000$ on all tested
encoders (Section~\ref{app:encoder_generalization}).

\textbf{Eighth}, a discrete-step SIR epidemic model
\citep{gu2024agentsmith} over a shared FAISS knowledge base shows that
a single poisoned entry amplifies via re-storage; \sad{}-based write-time
quarantine eliminates secondary infection for high-AUROC attack-defense pairs.

\textbf{Ninth}, we extend the threat model to graph-structured memory by
implementing three topology-exploiting attacks (hub insertion, edge-hijack,
subgraph cluster) and two graph-native defenses (degree-anomaly detection,
adjacency contamination scoring) over a GraphRAG-style entity-relation store.

\textbf{Tenth}, FPR is validated over 20 independent trials using Clopper-Pearson
exact binomial confidence intervals \citep{clopper1934use}, and production-scale
$\asra$ is directly measured using a GPT-4o-mini agent evaluator with
\texttt{text-embedding-3-large} retrieval.

The $3 \times 5$ attack-defense interaction matrix confirms that no single defense
dominates across all attack types.  Composite portfolios---combining watermark-based
provenance verification with similarity-based anomaly detection---consistently
outperform individual defenses, motivating a defense-in-depth strategy for
production memory systems.

\paragraph{Open challenges.}
One key challenge remains: \emph{formal robustness bounds for composite defenses}.
While the gradient-coupling property holds for \sad{} in isolation
(Proposition~\ref{prop:gradient_tension}), no formal guarantee exists for the
composite portfolio under adversaries that simultaneously optimize against all
defense components.  Extending the coupling proof to the composite objective
is a natural direction for future theoretical work.

\paragraph{Reproducibility.}
All code, the 200-entry synthetic corpus, evaluation scripts, and generated figures
are open-sourced at \url{https://github.com/ishrith-gowda/memory-agent-security}.
The repository includes a full pipeline script (\texttt{src/scripts/run\_pipeline.py})
that reproduces all tables and figures from a single command.

As autonomous agents increasingly mediate consequential decisions---scheduling,
information retrieval, task planning---memory integrity becomes a safety-critical
property.  We hope this framework provides a rigorous foundation for future research
on this underexplored attack surface.
